% Generated by ChimCanhCut Research Platform
% Template: IEEEtran — IEEE Transactions on Neural Networks and Learning Systems / TPAMI / TKDE
% Topic: Federated Learning under Concept Drift: Taxonomy, Challenges, and Future Directions

\documentclass[journal]{IEEEtran}

\usepackage[utf8]{inputenc}
\usepackage{cite}
\usepackage{amsmath,amssymb}
\usepackage{graphicx}
\usepackage{booktabs}
\usepackage{hyperref}
\usepackage{url}
\usepackage{multirow}
\usepackage{array}

\title{Federated Learning under Concept Drift: Taxonomy, Challenges, and Future Directions}
\author{[Author Names]}

\begin{document}
\maketitle

\begin{abstract}
% TODO: Enter your abstract via the Paper Metadata panel in the UI.
% Required: 150-250 words covering motivation, methodology, key findings, contributions.
\end{abstract}

\begin{IEEEkeywords}Federated Learning, Concept Drift, Non-IID Data, Distribution Shift\end{IEEEkeywords}

\section{Introduction}
\label{sec:introduction}

Federated Learning (FL) has emerged as a transformative paradigm for training machine learning models across distributed networks, enabling collaborative optimization without the need to centralize sensitive raw data [CITE:327, CITE:344, CITE:801, CITE:802]. This privacy-preserving architecture is particularly advantageous for edge computing and Internet of Things (IoT) environments, where data is generated locally on resource-constrained devices [CITE:451, CITE:592]. However, real-world data in these environments is rarely static; it typically arrives in continuous streams and evolves over time [CITE:195, CITE:393]. To navigate this complexity, it is crucial to establish a rigorous taxonomy of data variations early on, explicitly disambiguating spatial heterogeneity from temporal heterogeneity. Spatial heterogeneity, formally termed *Concept Skew* (or non-IID data), dictates that the joint data distribution varies across the network at any given time, such that $P_i(X, y) \neq P_j(X, y)$ for any two distinct clients $i$ and $j$ [CITE:217, CITE:584, CITE:806]. Conversely, temporal heterogeneity is broadly defined as *Concept Drift*, representing a change in the joint probability distribution of the input data $X$ and the target variable $y$ over time for a given client. Mathematically, general concept drift occurs when $P_t(X, y) \neq P_{t+\Delta t}(X, y)$ for a given time step $t$ and interval $\Delta t$ [CITE:252, CITE:328, CITE:803]. This temporal non-stationarity can be further decomposed into two primary manifestations: *covariate shift*, where the input data distribution changes ($P_t(X) \neq P_{t+\Delta t}(X)$) while the conditional distribution remains static; and *concept shift* (often termed real concept drift or label shift), where the underlying decision boundary itself evolves, meaning $P_t(y|X) \neq P_{t+\Delta t}(y|X)$ [CITE:804, CITE:805]. 

The intersection of FL and temporal concept drift introduces profound complexities that challenge the foundational assumptions of distributed optimization. In centralized machine learning, concept drift is typically managed using statistical process control algorithms like the Drift Detection Method (DDM) or Early Drift Detection Method (EDDM), which rely on continuous, real-time monitoring of raw error rates or data streams [CITE:830, CITE:831]. However, applying these centralized paradigms directly to FL environments is fundamentally incompatible due to the unique architectural constraints of distributed networks. First, strict privacy requirements prohibit the central server from inspecting raw client data or localized error streams, effectively blinding the server to the onset of localized concept shifts. Second, FL systems typically operate under partial client participation, meaning the server only observes a fraction of the network during any given communication round; this leads to stale client models and masks underlying drift dynamics. Finally, communication bottlenecks introduce significant latency, delaying the aggregation of drift signals and rendering reactive adaptations obsolete by the time they are globally deployed [CITE:832, CITE:833]. Consequently, spatiotemporal drift in FL requires fundamentally new detection and adaptation mechanisms.

To mitigate these spatiotemporal challenges, the literature has proposed various adaptation strategies. Traditional FL frameworks rely heavily on synchronous parameter aggregation \cite{Li2019}. While early literature suggested asynchronous updates were strictly necessary for non-stationary streaming data, recent advancements demonstrate that synchronous FL frameworks equipped with straggler mitigation, proximal regularization terms \cite{paper826}, or normalized updates \cite{paper827} remain highly viable. Nevertheless, to bypass latency bottlenecks in highly dynamic edge networks, a parallel body of literature advocates for continuous, semi-synchronous, or fully asynchronous learning frameworks [CITE:155, CITE:174, CITE:592, CITE:810]. Other approaches leverage explicit unsupervised and partially-supervised drift detection algorithms [CITE:330, CITE:541, CITE:811], or Clustered Federated Learning (CFL) to group participants with similar data distributions, thereby training specialized models resilient to localized drifts [CITE:216, CITE:218, CITE:812]. Furthermore, robust representation learning frequently utilizes lightweight, sequential algorithms tailored for resource-constrained devices [CITE:349, CITE:607, CITE:813]. Interestingly, while most frameworks rely on sharing model parameters or gradients, some recent approaches explicitly avoid parameter sharing to mitigate label heterogeneity. Instead, they share alternative representations such as "diagnostic knowledge," which practically involves transmitting aggregated, class-conditional feature statistics or soft-label distributions, allowing the server to model global data distribution properties without exposing raw samples [CITE:503, CITE:607, CITE:814].

Despite these methodological advancements, a critical synthesis of the current literature reveals significant contradictions, particularly regarding the disconnect between theoretical models of concept drift and practical FL implementations. Theoretical frameworks, such as those grounded in statistical mechanics \cite{Straat2018} and extended in recent analyses [CITE:815, CITE:816], warn of persistent sub-optimal plateau states caused by continuous drift. This theory-practice gap exists primarily because theoretical bounds often rely on strict thermodynamic limits, infinite sample sizes, or purely Markovian online learning assumptions where past data is permanently lost, leading to inevitable catastrophic forgetting \cite{paper820}. Conversely, applied frameworks frequently report continuous performance improvements and the successful mitigation of catastrophic forgetting through practical workarounds, such as the deployment of global replay buffers augmented with synthetic data generation [CITE:195, CITE:817]. By artificially re-injecting historical data distributions into the current training step, these replay mechanisms effectively break the strict Markovian assumptions of the theoretical models, converting a non-stationary online problem into a pseudo-stationary one. However, these generative global replay buffers introduce their own critical vulnerabilities; they are highly susceptible to bias amplification and quality degradation, which are strictly bounded by the generative capabilities of the global model [CITE:818, CITE:819]. 

Empirical evaluation in this domain also suffers from severe limitations. There is a critical absence of unified, large-scale, real-world streaming datasets specifically designed to benchmark temporal and distributed concept drift in FL [CITE:821, CITE:822]. Current evaluations rely heavily on highly fragmented, domain-specific, or entirely synthetic datasets, making cross-study comparisons highly challenging. Additionally, multi-model and clustering adaptation approaches often rely on unrealistic assumptions, such as having oracle knowledge of the exact number of concepts (clusters) or requiring the communication of balanced classifiers, which are rarely available in dynamic environments [CITE:140, CITE:141, CITE:823]. Crucially, these approaches frequently struggle with severe "cold-start" performance degradation when new concepts emerge, as newly created specialized models lack sufficient historical training data to generalize effectively [CITE:824, CITE:825].

To the best of our knowledge, existing surveys predominantly focus either on spatial data heterogeneity (concept skew) in static FL environments [CITE:840, CITE:841] or on concept drift in purely centralized machine learning settings [CITE:842, CITE:843]. While a few recent reviews touch upon continual federated learning, they primarily address task-incremental learning rather than continuous, unsegmented streaming data subject to sudden or gradual concept shifts [CITE:844, CITE:845]. To bridge this gap, this paper presents the first comprehensive survey dedicated specifically to Federated Learning under concept drift. Our primary contributions are threefold: (1) We identify and formalize the unique challenges imposed by FL constraints (e.g., privacy, partial participation, latency) on traditional drift detection methodologies; (2) We provide a critical analysis of the theoretical-practical divide in federated drift adaptation, highlighting the limitations of current generative replay and clustering methods; and (3) We propose a unified, multi-dimensional taxonomy to systematically categorize the state-of-the-art.

To address fragmented research efforts and guide future exploration, our proposed taxonomy categorizes the literature across four primary dimensions, which directly mirror the structural organization of the subsequent sections in this paper. Section 2 explores **Learning Architectures**, differentiating between synchronous, semi-synchronous, asynchronous, and decentralized topologies. Section 3 examines **Adaptation Strategies**, classifying methods into active drift detection, passive continuous updates, clustering, and memory-based replay mechanisms. Section 4 discusses **Application Domains**, highlighting domain-specific constraints in IoT, healthcare, and autonomous systems. Section 5 analyzes **Data Characteristics**, detailing the severity, speed, and spatial distribution of the drift. Finally, Section 6 concludes the survey by outlining open challenges—such as cold-start degradation, generative replay vulnerabilities, and the need for standardized benchmarking—providing a foundational roadmap for developing robust, drift-resilient federated learning systems.

\begin{figure}[!t]
\centering
\includegraphics[width=\columnwidth]{figures/year_distribution}
\caption{Publication year distribution of included papers.}
\label{fig:year_dist}
\end{figure}

\section{Background and Related Work}
\label{sec:background}

**Background**

Federated Learning (FL), introduced by McMahan et al. \cite{Kairouz2020}, has become a prominent privacy-preserving paradigm, enabling distributed clients to collaboratively train machine learning models without exchanging raw local data. Conventional FL architectures are typically classified into horizontal, vertical, and transfer learning, depending on how feature spaces and sample identifiers align across participants \cite{NERGZ2022}, \cite{paper801}, \cite{paper802}. As FL deployment expands into edge computing and the Internet of Things (IoT), it increasingly relies on decentralized, cross-device architectures to integrate machine learning with sophisticated wireless communication strategies \cite{Letaief2021}, \cite{paper803}. However, standard aggregation algorithms, such as Federated Averaging (FedAvg), implicitly assume static data distributions—an assumption that often fails in dynamic, real-world environments where data streams are inherently non-stationary.

A fundamental challenge in FL is spatial data heterogeneity, widely known as non-IID (Independent and Identically Distributed) data \cite{Deshwal2025}. Because clients operate in diverse environments, their local datasets exhibit significant statistical variation. Formally, let $\mathcal{D}_i = \{(x_k, y_k)\}_{k=1}^{n_i}$ represent the local dataset of client $i$, drawn from a joint probability distribution $P_i(X, Y)$. Spatial heterogeneity implies that for any two clients $i$ and $j$ ($i \neq j$), the joint distributions differ: $P_i(X, Y) \neq P_j(X, Y)$ \cite{Gutierrez2025}, \cite{paper804}. Applying Bayes' theorem, this joint distribution can be decomposed as $P(X, Y) = P(Y|X)P(X) = P(X|Y)P(Y)$. Consequently, spatial heterogeneity manifests as feature skew ($P_i(X) \neq P_j(X)$), label skew ($P_i(Y) \neq P_j(Y)$), or concept skew ($P_i(Y|X) \neq P_j(Y|X)$). To prevent terminological confusion, it is imperative to explicitly disambiguate spatial *concept skew* from temporal *concept shift* at this juncture. Concept skew denotes spatial heterogeneity, where the underlying decision boundary $P(Y|X)$ differs across distinct clients at the exact same point in time. In contrast, concept shift refers exclusively to temporal heterogeneity, where a single client's or the global model's $P(Y|X)$ evolves over time. To mitigate spatial heterogeneity, researchers have proposed robust aggregation methods such as FedProx \cite{Cao2025} and SCAFFOLD \cite{Karimireddy2019a}, the latter of which employs control variates to refine local updates and reduce variance. Other notable approaches include Personalized Federated Learning (PFL) \cite{Tan2022}, \cite{paper805} and Clustered Federated Learning (CFL) \cite{Liu2025}, \cite{Yan2023a}. While these methods effectively address static non-IID data, they do not account for temporal distribution shifts.

Temporal heterogeneity, broadly defined as *concept drift*, occurs when the statistical properties of the target variable or input data evolve over time within streaming environments \cite{Mannapur2025}, \cite{Agrahari2021}. Mathematically, concept drift is defined as a change in the joint probability distribution over time: $P_t(X, Y) \neq P_{t+\Delta t}(X, Y)$ for a given time step $t$ and interval $\Delta t$ \cite{Hinder2023}, \cite{Mannapur2025}, \cite{paper803}. Based on the decomposition of this joint distribution, concept drift is formally categorized into three primary sub-types \cite{paper806}, \cite{paper807}:
1. **Covariate Shift**: The input feature distribution changes over time ($P_t(X) \neq P_{t+1}(X)$), while the conditional distribution of labels given features remains stable ($P_t(Y|X) = P_{t+1}(Y|X)$).
2. **Label Shift** (Prior Probability Shift): The distribution of the target variable changes ($P_t(Y) \neq P_{t+1}(Y)$), but the conditional distribution of features given labels remains constant ($P_t(X|Y) = P_{t+1}(X|Y)$).
3. **Concept Shift** (Real Drift): The fundamental mapping or decision boundary between features and labels changes over time ($P_t(Y|X) \neq P_{t+1}(Y|X)$). 

It is critical to maintain theoretical rigor regarding these definitions: *concept drift* serves as the umbrella term for any temporal change in the joint distribution $P(X,Y)$, whereas *concept shift* strictly denotes a change in the conditional distribution $P(Y|X)$. 

Concept drift manifests in various temporal forms, including abrupt shifts, gradual trends, and recurring concepts \cite{Prasad2024}. In centralized machine learning, Gama et al. \cite{Hinder2024} and Lu et al. \cite{paper808} observe that drift is traditionally managed through explicit detection methodologies. These typically involve a four-staged scheme: data windowing, descriptor building, test statistic calculation, and adaptation. Prominent examples include the Drift Detection Method (DDM) and the Early Drift Detection Method (EDDM) \cite{paper809}, which continuously monitor the raw error rates of a model on a centralized data stream to trigger retraining when statistical thresholds are breached.

While these centralized techniques are highly effective when data is aggregated in a single repository, transitioning them to a federated environment introduces profound architectural friction. Centralized methods implicitly assume continuous, unrestricted access to the global data stream and the ability to evaluate model performance instantaneously across all available samples. In Federated Learning, however, the data is fundamentally decentralized, siloed across heterogeneous edge devices, and shielded by strict privacy protocols. This paradigm shift invalidates the core assumptions of traditional drift detection, necessitating a fundamental redesign of how non-stationarity is identified and managed.

Specifically, directly applying traditional drift detection algorithms to FL is fundamentally impractical due to three unique spatiotemporal challenges and architectural constraints. First, **privacy constraints** strictly prohibit the central server from inspecting raw data streams, invalidating the centralized data windowing or raw error-rate tracking utilized by DDM and EDDM. Second, **partial client participation** dictates that the server only aggregates updates from a small, often non-representative subset of clients per round, leading to biased or severely delayed global drift detection. Third, **communication bottlenecks** restrict the frequent exchange of diagnostic knowledge. In a practical FL setting, sharing diagnostic knowledge requires transmitting localized loss landscapes, gradient variance metrics, or explicit drift alerts (e.g., binary drift flags) to the server; the high latency and bandwidth costs of transmitting these metrics hinder real-time responsiveness \cite{paper810}. Consequently, FL systems require specialized mechanisms to detect and adapt to the simultaneous evolution of the joint data distribution across distributed nodes.

Despite the critical nature of this problem, existing literature lacks a comprehensive review at the intersection of FL and concept drift. To the best of our knowledge, existing surveys predominantly focus either on handling static non-IID data in FL \cite{Deshwal2025}, \cite{paper804} or on concept drift within centralized and traditional distributed machine learning paradigms \cite{Hinder2024}, \cite{paper808}. While a few recent studies touch upon dynamic environments in FL \cite{paper811}, they do not provide a systematic categorization of drift adaptation mechanisms under strict federated constraints. This paper fills this critical gap by contributing a comprehensive taxonomy of FL under concept drift and detailing how federated constraints necessitate novel algorithmic designs.

Specifically, we structure our survey around a novel multi-dimensional taxonomy encompassing four key dimensions that precisely align with the subsequent core body sections of this paper: (1) **Learning Architectures**, detailing the structural paradigms (e.g., horizontal, vertical, decentralized) utilized to manage dynamic data streams; (2) **Adaptation Strategies**, exploring how server-side aggregation and client-side local models dynamically adjust to temporal shifts without violating privacy; (3) **Application Domains**, highlighting the real-world edge and IoT environments where federated concept drift is most prevalent; and (4) **Data Characteristics**, summarizing the specific types of drift, spatiotemporal heterogeneities, and evaluation datasets required to rigorously assess dynamic FL systems. By systematically exploring these four dimensions, this survey aims to elucidate current challenges and chart future directions for robust, drift-resilient federated learning.

\section{Problem Formulation}
\label{sec:problem_formulation}

The problem of Federated Learning (FL) under concept drift fundamentally extends the traditional distributed optimization objective by introducing spatiotemporal non-stationarity. In standard FL, the goal is to learn a global model parameter $w$ by minimizing an empirical risk function over data distributed across $N$ clients. Let $\mathcal{X} \subset \mathbb{R}^d$ denote the feature space and $\mathcal{Y}$ denote the label space. The objective is formally defined as $\min_w F(w) = \sum_{i=1}^N p_i \mathbb{E}_{(x,y)\sim \mathcal{D}_i} [\ell(w; x, y)]$, where $p_i$ is the weight of the $i$-th client, $\ell$ is the loss function, and $\mathcal{D}_i$ is the local data distribution. Traditional FL inherently deals with spatial heterogeneity—formalized as non-Independent and Identically Distributed (non-IID) data across clients. To ensure theoretical rigor and prevent terminology collisions, this spatial heterogeneity must be explicitly disambiguated from temporal changes early in our formulation. Spatial non-stationarity often manifests as *Concept Skew*, where the conditional distribution of labels given features differs across clients ($P_i(Y|X) \neq P_j(Y|X)$ for clients $i \neq j$) even if the marginal feature distributions remain similar [CITE:217, 441]. However, standard FL formulations operate under the strict assumption that these local distributions, regardless of their spatial skew, remain stationary over time. In real-world streaming environments, such as edge computing and Internet of Things (IoT) networks, this stationary assumption is frequently violated [CITE:328, 349].

General *Concept Drift* introduces temporal heterogeneity, occurring when the joint probability distribution of features and labels changes over time. Mathematically, this overarching phenomenon is expressed as $P_{i,t}(X, Y) \neq P_{i,t+1}(X, Y)$ for a given client $i$ at time $t$ \cite{Hinder2024}. This temporal non-stationarity must be formally decomposed into distinct dynamics: covariate shift, where the input distribution changes ($P_{i,t}(X) \neq P_{i,t+1}(X)$) while the conditional distribution remains stable; prior probability shift ($P_{i,t}(Y) \neq P_{i,t+1}(Y)$); and true *Concept Shift*, where the underlying decision boundary itself evolves over time ($P_{i,t}(Y|X) \neq P_{i,t+1}(Y|X)$). While spatial Concept Skew dictates how models generalize across different clients at a static point in time, temporal Concept Shift dictates how models adapt to evolving environments. These temporal shifts can manifest as abrupt changes, gradual trends, or recurring concepts [CITE:135, 296, 559]. Consequently, the learning objective must transition from minimizing a static empirical risk to dynamically minimizing the risk over a continuous, evolving stream of data: $\min_{w_t} \sum_{i=1}^N p_{i,t} \mathbb{E}_{(x,y)\sim \mathcal{D}_{i,t}} [\ell(w_t; x, y)]$ [CITE:195, 393].

A core challenge in formulating this dynamic objective is managing the intersection of spatial and temporal heterogeneity, particularly regarding drift detection. In traditional, centralized machine learning, concept drift is typically managed using established statistical detectors like the Drift Detection Method (DDM) or Early Drift Detection Method (EDDM) [CITE:112, 284]. These centralized approaches rely on continuous, unrestricted access to raw data streams and global error rates to monitor distribution changes. 

Transitioning these centralized detection mechanisms to a federated environment, however, introduces severe structural and privacy constraints. Privacy mandates and communication bottlenecks strictly prohibit the sharing of raw streaming data, rendering centralized DDM and EDDM inapplicable in their native forms [CITE:325, 410]. Consequently, FL dictates that changes in the joint data distribution must be inferred indirectly. Systems must instead rely on localized drift detection at the client level or the privacy-preserving aggregation of drift signals—such as model weight divergences or localized error statistics—at the server [CITE:488, 512].

This dynamic formulation complicates traditional round-based FL aggregation and brings the stability-plasticity dilemma to the forefront [CITE:154, 318]. When a new concept emerges, models often suffer from a "cold-start" performance drop due to a lack of training on the new distribution. Conversely, rapid adaptation to new concepts can lead to catastrophic forgetting of previously learned knowledge. To manage this, the choice of synchronization strategy is critical. While asynchronous updates are highly beneficial for handling continuous streaming data and mitigating stragglers in edge environments without waiting for global synchronization barriers [CITE:196, 592], many synchronous FL frameworks successfully operate in drifting environments by incorporating advanced proximal regularization terms (e.g., variants of FedProx or FedNova) [CITE:412, 501]. These proximal terms constrain local updates to prevent the model from overfitting to transient, localized drifts, thereby providing a balanced approach to the stability-plasticity trade-off.

Critically, analyzing this domain reveals a significant disconnect between theoretical formulations of concept drift in federated settings and practical implementations. Theoretical analyses and generalization bounds for FL under concept drift currently rely on strict statistical assumptions—such as infinite sample sizes, sub-Gaussian loss functions, or bounded drift rates ($\sum_{t} ||\mathcal{D}_t - \mathcal{D}_{t-1}|| \le \epsilon$)—which are rarely applicable to complex, finite-sample real-world data streams. For instance, theoretical frameworks utilizing statistical mechanics warn of persistent, sub-optimal "plateau states" caused by concept drift \cite{Straat2018}. This theoretical limitation arises because analytical models often assume strict Markovian drift, where the model constantly chases a moving target distribution without any historical memory, leading to an unavoidable lower bound on the generalization error. 

In contrast, applied frameworks frequently bypass these theoretical plateau states by integrating continuous learning mechanisms that explicitly break the assumption of strict Markovian drift. Practical formulations reintroduce historical context through global replay buffers [CITE:195, 310] or dynamic repositories of specialized models \cite{Tsiporkova2023}. By storing subsets of historical data ($\mathcal{D}_{i, t-k}$), these empirical methods effectively smooth the temporal distribution, transforming a strictly non-stationary problem into a locally stationary one over a sliding window. Furthermore, to accommodate the computational constraints of edge devices during rapid concept shifts, empirical solutions increasingly adopt lightweight sequential learning architectures that bypass iterative backpropagation [CITE:349, 543]. To prevent negative transfer during aggregation, recent frameworks advocate for sharing generalized "diagnostic knowledge" rather than raw parameters. This includes transmitting localized error statistics, feature importance metrics, or intermediate representations, allowing the global server to identify and aggregate only the parameters relevant to the current active concept [CITE:607, 612].

Addressing these theoretical gaps while accommodating the realities of asynchronous and synchronous non-IID streaming data remains the central focus of problem formulation in this domain. To systematically navigate these intersecting challenges—specifically, the dual impact of spatial Concept Skew and overarching temporal Concept Drift, the constraints of decentralized drift detection, and the persistent theory-practice divide—a structured categorization of existing literature is required. Consequently, the subsequent sections of this survey propose and detail a multi-dimensional taxonomy encompassing four key dimensions: Learning Architectures, Adaptation Strategies, Application Domains, and Data Characteristics. This taxonomy is designed to explicitly classify and evaluate current methodologies in Federated Learning under concept drift along these critical axes.

\section{Methodology}
\label{sec:methodology}

\# Methodology

This systematic survey adheres to the Preferred Reporting Items for Systematic Reviews and Meta-Analyses (PRISMA) 2020 guidelines, ensuring transparency, reproducibility, and rigor in synthesizing literature on Federated Learning (FL) under concept drift.

\#\# 1. Search Strategy
To capture the evolving landscape of FL and non-stationary data distributions, we conducted a comprehensive search across Semantic Scholar, arXiv, and IEEE Xplore. Our strategy combined controlled vocabulary with free-text keywords, including "Federated Learning," "Concept Drift," "Non-IID data," "Online Learning," and "Model Adaptation." 

We constructed search queries to balance precision and recall, employing Boolean operators (AND/OR) to account for terminological variations across machine learning and distributed systems. The search period spans January 2015 to January 2026, reflecting the emergence of FL as a distinct research field following the seminal work of McMahan et al. [CITE:foundational_01].

\#\# 2. PRISMA Flow
The selection process followed the four-stage PRISMA flow diagram illustrated in Figure 1:

*   **Identification:** Database searches yielded 490 records.
*   **Screening:** After removing duplicates, 490 records underwent initial screening. We excluded 78 records via rule-based filtering (e.g., non-English language, pre-2015 publication date) and 71 records through automated, LLM-assisted screening, which identified papers lacking substantive relevance to the intersection of FL and concept drift.
*   **Eligibility:** We assessed the remaining 341 records for full-text eligibility. All 341 papers met the inclusion criteria, providing sufficient technical depth regarding distributed learning in non-stationary environments.
*   **Included:** The final corpus comprises 341 studies, categorized into three tiers: 84 studies directly addressing FL under concept drift, 250 studies providing adjacent contributions (e.g., drift detection in centralized settings or FL under non-IID data without explicit drift), and 7 foundational papers establishing the field's theoretical underpinnings.

\#\# 3. Inclusion and Exclusion Criteria
To maintain the survey's integrity, we applied the following criteria:

**Inclusion Criteria:**
*   **Peer-reviewed or arXiv preprints:** To capture both established methodologies and cutting-edge developments.
*   **Relevance:** Focus on papers addressing FL, concept drift, or related topics such as catastrophic forgetting and online model updates.
*   **Temporal Scope:** Publications between 2015 and 2026 to ensure alignment with the rapid evolution of FL architectures.
*   **Accessibility:** Availability of an abstract and full text for comprehensive analysis.

**Exclusion Criteria:**
*   **Duplicate entries:** To prevent synthesis bias.
*   **Non-English papers:** To ensure linguistic consistency and accurate technical interpretation.
*   **Pre-2015 publications:** To focus on the modern era of FL research.
*   **Inaccessible abstracts:** Papers lacking sufficient metadata were excluded to avoid incomplete data extraction.

\#\# 4. Quality Assessment
Given the diversity of publication venues—ranging from top-tier conferences like NeurIPS and ICML to specialized workshops—we implemented a multi-tiered quality assessment. Each paper was evaluated based on: (1) **Methodological Rigor:** The clarity of the proposed FL architecture and the robustness of the drift detection mechanism; (2) **Experimental Validity:** The use of standard benchmarks (e.g., FEMNIST, CIFAR-10, or synthetic drift datasets) and the presence of comparative baselines; and (3) **Contribution Significance:** The extent to which the paper addresses the probability distribution of the input data $X$ and the target variable $y$. 

Mathematically, temporal heterogeneity occurs when $P_t(X, y) \neq P_{t+\Delta t}(X, y)$ for a given time step $t$ and interval $\Delta t$ [CITE:252, CITE:328, CITE:803]. This non-stationarity manifests as covariate shift, where the input data distribution changes ($P_t(X) \neq P_{t+\Delta t}(X)$), or as posterior drift, where the underlying decision boundary evolves ($P_t(y|X) \neq P_{t+\Delta t}(y|X)$).

\begin{figure}[!t]
\centering
\includegraphics[width=0.85\columnwidth]{figures/prisma_flow}
\caption{PRISMA 2020 flow diagram of the systematic literature search.}
\label{fig:prisma}
\end{figure}

\section{Taxonomy}
\label{sec:taxonomy}

To systematically analyze the intersection of Federated Learning (FL) and non-stationary environments, it is imperative to establish a formal mathematical foundation that rigorously distinguishes between spatial and temporal data heterogeneity. Let $\mathcal{X} \in \mathbb{R}^d$ denote the feature space and $\mathcal{Y}$ denote the label space. 

In standard FL, spatial heterogeneity (non-IID data) occurs when the joint data distributions across clients $i$ and $j$ diverge, such that $P_i(X, Y) \neq P_j(X, Y)$. This spatial skew can be further decomposed into feature skew ($P_i(X) \neq P_j(X)$), label skew ($P_i(Y) \neq P_j(Y)$), and spatial **Concept Skew**, where the conditional distribution of labels given features varies across clients at a static point in time ($P_i(Y|X) \neq P_j(Y|X)$). 

When introducing temporal dynamics, the system must also contend with temporal heterogeneity, formally defined as **Concept Drift**. Concept drift occurs when the joint distribution at a given client (or globally) changes over time, meaning $P_t(X, Y) \neq P_{t+1}(X, Y)$. To prevent terminology collisions prevalent in the literature, it is crucial to explicitly disambiguate spatial *Concept Skew* from temporal *Concept Shift* early in the architectural design. While Concept Skew refers to conditional discrepancies across different clients simultaneously, Concept Drift encompasses three distinct temporal phenomena:
1.  **Covariate Shift:** The distribution of features changes over time ($P_t(X) \neq P_{t+1}(X)$), but the conditional distribution of the labels remains constant ($P_t(Y|X) = P_{t+1}(Y|X)$).
2.  **Prior Probability (Label) Shift:** The distribution of labels changes ($P_t(Y) \neq P_{t+1}(Y)$), while the conditional distribution of features given the labels remains constant ($P_t(X|Y) = P_{t+1}(X|Y)$).
3.  **Temporal Concept Shift (True Concept Drift):** The fundamental relationship between features and labels changes temporally, meaning the conditional distribution shifts over time ($P_t(Y|X) \neq P_{t+1}(Y|X)$).

Because existing literature often conflates spatial Concept Skew with temporal Concept Drift, or addresses specific drift types in isolation without acknowledging their interplay, a structured classification is essential. We propose a multi-dimensional taxonomy to provide a unified framework that bridges these isolated approaches. This taxonomy categorizes the state-of-the-art across four primary dimensions: Learning Architecture, Adaptation Strategy, Data Characteristics, and Application Domain [CITE:131, CITE:279].

**1. Learning Architecture**
The structural topology of the FL system fundamentally dictates how temporal concept drift is managed and mitigated.
*   **Edge and Cross-Device FL:** Frequently deployed in resource-constrained Internet of Things (IoT) environments, this architecture demands communication-efficient and computationally lightweight model updates [CITE:159, CITE:349]. Traditional synchronous parameter averaging struggles to keep pace with real-time non-stationary environments; thus, asynchronous online federated learning (e.g., ASO-Fed) is highly beneficial for handling continuous streaming data and mitigating the impact of stragglers \cite{Chen2020}. However, asynchronous updates are not strictly necessary to handle drift. Synchronous FL frameworks equipped with straggler mitigation, proximal regularization, or normalized updates (e.g., variants of FedProx or FedNova) remain actively researched. By bounding local updates and penalizing extreme deviations, these synchronous methods successfully stabilize training and mitigate client drift even in highly dynamic environments [CITE:801, CITE:802].
*   **Cross-Silo FL:** Involves collaboration among distinct organizations, such as medical institutions, where data distributions may shift over time due to changing demographics, updated clinical protocols, or new equipment [CITE:382, CITE:512].
*   **Decentralized and Clustered FL:** Clustered Federated Learning (CFL) groups participants with similar data distributions to train specialized models. This effectively localizes the impact of drift and improves personalization by isolating shifting concepts to specific client clusters, preventing global model degradation [CITE:216, CITE:222, CITE:491].

**2. Adaptation Strategy**
Approaches to detecting and mitigating concept drift vary significantly across the literature, ranging from explicit detection to continuous adaptation.
*   **Explicit Drift Detection:** Historically, methodologies have relied on lightweight on-device detectors like ADWIN, DDM, or EDDM to monitor abrupt shifts and gradual trends [CITE:135, CITE:138, CITE:328]. 
    *Transitioning these centralized detection mechanisms to a federated setting, however, introduces significant algorithmic friction.* Centralized detectors inherently assume uninhibited access to raw, continuous data streams and global error rates. In an FL environment, relying on such assumptions violates core privacy constraints and incurs prohibitive communication overhead. Consequently, modern FL frameworks must adapt these methods by utilizing privacy-preserving federated statistical tests or leveraging aggregated Shapley values to continuously evaluate shifting client data valuation without exposing raw features [CITE:412, CITE:549, CITE:610].
*   **Online and Continuous Model Update:** Because static, round-based FL is inadequate for streaming data, frameworks increasingly adopt continuous learning mechanisms. Applied frameworks successfully mitigate catastrophic forgetting using practical techniques like global replay buffers, synthetic data generation, and Online Sequential Extreme Learning Machines (OS-ELM). OS-ELM is particularly valuable in this context; it provides lightweight, sequential updating capabilities that allow models to adapt rapidly to new data without the computational overhead of extensive backpropagation—a critical advantage for resource-constrained FL nodes experiencing drift [CITE:195, CITE:349, CITE:418].
*   **Personalization \& Clustering:** Addressing label heterogeneity is frequently treated as a parallel challenge to concept drift. Solutions include dynamic client clustering [CITE:130, CITE:458], fine-grained layer aggregation \cite{Fu2025}, and hierarchical aggregation paradigms \cite{Aziz2025}.
*   **Robust Representation Learning:** Focuses on extracting invariant features to withstand drift. Autoencoder architectures are repeatedly utilized for anomaly detection and robust representation under drifting conditions [CITE:349, CITE:607]. Other methods avoid standard parameter sharing to mitigate negative transfer, opting instead to share "diagnostic knowledge." This alternative representation entails transmitting abstracted meta-information—such as class-conditional feature distributions, prototype vectors, or soft labels—rather than raw model weights. By doing so, these methods decouple global feature alignment from local classifier optimization, preserving local drift adaptations while still benefiting from global insights [CITE:503, CITE:518, CITE:607].
*   **Theoretical Analysis and the Theory-Practice Gap:** Theoretical literature provides foundational bounds on learning dynamics, utilizing information-theoretic approaches to balance empirical risk minimization and generalization under drift [CITE:132, CITE:590]. However, a critical theory-practice gap exists regarding catastrophic forgetting. Statistical mechanics frameworks often model online learning under strict Markovian assumptions (i.e., memoryless updates) using idealized data-generating processes. Under these strict theoretical constraints, models are mathematically proven to become trapped in persistent sub-optimal plateau states \cite{Straat2018}. Conversely, applied frameworks frequently bypass these pessimistic bounds by intentionally violating strict memoryless assumptions. For example, the empirical success of global replay buffers breaks the Markovian assumption of the data stream, artificially restoring a degree of stationarity to the training process and effectively preventing predicted plateau states [CITE:195, CITE:349].

**3. Data Characteristics**
The mathematical nature of the data heavily influences the algorithmic design of drift-aware FL systems.
*   **Continuous Streaming Data:** Requires real-time processing and continuous model adaptation, often handled via asynchronous updates and online learning algorithms designed for environments where the joint distribution $P_t(X, Y)$ evolves rapidly and unpredictably [CITE:196, CITE:282].
*   **Static Non-IID Data:** Spatial heterogeneity ($P_i(X, Y) \neq P_j(X, Y)$) is frequently addressed alongside temporal heterogeneity ($P_t(X, Y) \neq P_{t+1}(X, Y)$). This dual challenge requires robust aggregation methods and classifier calibration to disentangle local data skew from genuine global concept drift [CITE:571, CITE:584].
*   **Spatio-Temporal Data:** Common in traffic analysis and sensor networks, necessitating specialized architectures like spatio-temporal memory-enhanced autoencoders or federated graph neural networks to capture complex, multi-dimensional shifts [CITE:195, CITE:375].
*   **Partially Labeled and Imbalanced Data:** Drift detection and adaptation must often operate under limited labeling or multi-class imbalance, requiring semi-supervised learning or specialized skew-insensitive loss functions to prevent false drift alarms triggered by missing classes rather than true distributional shifts [CITE:294, CITE:298, CITE:553].

**4. Application Domain**
Federated Learning under concept drift is critical across several highly dynamic domains where data distributions are inherently non-stationary.
*   **Industrial IoT \& Sensor Networks:** Applications include collaborative fault diagnosis for wind turbines \cite{Jiang2024}, robust photovoltaic (PV) power forecasting under weather-induced drift \cite{Moayyed2025}, and distributed intrusion detection systems adapting to novel cyber threats [CITE:160, CITE:339].
*   **Healthcare \& Medical Imaging:** Managing evolving medical datasets and domain shifts across hospitals using self-supervised learning and domain adaptation to maintain diagnostic accuracy despite changing clinical environments [CITE:202, CITE:321, CITE:597].
*   **Web \& Recommender Systems:** Adapting to dynamic, diversified user preferences over time through hybrid FL and drift detection approaches, ensuring recommendations remain relevant as user behavior evolves [CITE:139, CITE:494].

**Critical Analysis of the Taxonomy**
A critical analysis of the current literature dictates that the joint data distribution $P(X, Y)$ cannot be treated as a static entity in real-world FL deployments. The proposed taxonomy reveals that while significant strides have been made in isolated dimensions—such as optimizing edge architectures or developing robust representations—there remains a distinct lack of holistic frameworks that address the full spectrum of non-stationarity. Specifically, the interplay between spatial Concept Skew and temporal Concept Drift requires more rigorous algorithmic disentanglement; systems currently struggle to differentiate between a client possessing a uniquely skewed static dataset and a client experiencing a genuine temporal shift. Furthermore, future research must bridge the identified theory-practice gap by developing mathematical models that accurately reflect the non-Markovian adaptation strategies currently driving empirical success in the field. Moving forward, the development of unified protocols that concurrently manage cross-client spatial skew and cross-time temporal drift, without violating strict privacy and communication constraints, will be paramount to the maturation of Federated Learning in dynamic environments.

\begin{table}[!t]
\renewcommand{\arraystretch}{1.3}
\caption{Proposed Multi-Dimensional Taxonomy}
\label{tab:taxonomy}
\centering
\begin{tabular}{p{3cm}p{9cm}}
\toprule
\textbf{Dimension} & \textbf{Categories} \\
\midrule
\textbf{Learning\_Architecture} & Edge/Cross-Device FL, Cross-Silo FL, Decentralized/Clustered FL, Non-Federated/Traditional Online Learning \\
\textbf{Adaptation\_Strategy} & Explicit Drift Detection, Online/Continuous Model Update, Personalization \& Clustering, Robust Representation Learning, Theoretical \& Statistical Analysis \\
\textbf{Application\_Domain} & Industrial IoT \& Sensor Networks, Healthcare \& Medical Imaging, Web \& Recommender Systems, General/Theoretical Frameworks \\
\textbf{Data\_Characteristics} & Continuous Streaming Data, Static Non-IID Data, Spatio-Temporal Data, Partially Labeled/Delayed Label Data \\
\bottomrule
\end{tabular}
\end{table}

\begin{figure}[!t]
\centering
\includegraphics[width=\columnwidth]{figures/taxonomy_heatmap}
\caption{Heatmap of paper distribution across taxonomy dimensions.}
\label{fig:taxonomy_heatmap}
\end{figure}

\section{Literature Review}
\label{sec:literature_review}

The literature surrounding Federated Learning (FL) in non-stationary environments has expanded significantly, driven by the necessity to address evolving data streams. To rigorously analyze this domain and prevent the terminology collisions prevalent in early literature, it is essential to establish formal mathematical definitions that explicitly disambiguate spatial and temporal data dynamics before examining current methodologies. 

In standard FL, spatial heterogeneity (non-IID data) occurs when the joint distribution of features $X$ and labels $Y$ differs across clients, such that $P_i(X, Y) \neq P_j(X, Y)$ for clients $i \neq j$. A particularly challenging subset of this is **Spatial Concept Skew**, where the conditional distribution varies across the network ($P_i(Y|X) \neq P_j(Y|X)$). When continuous streaming data is introduced, systems must also contend with temporal heterogeneity, broadly defined as **General Concept Drift**. Formally, this occurs when the joint distribution for a given client $i$ evolves over time $t$, such that $P_{i, t}(X, Y) \neq P_{i, t+\Delta t}(X, Y)$. Based on Bayesian decision theory, where $P(X, Y) = P(Y|X)P(X)$, this temporal evolution manifests in three primary forms: *covariate shift*, where the input distribution changes ($P_{i,t}(X) \neq P_{i,t+\Delta t}(X)$) but the conditional distribution remains stable; *prior probability shift*, where the label distribution changes ($P_{i,t}(Y) \neq P_{i,t+\Delta t}(Y)$); and **Temporal Concept Shift** (often termed true concept drift), where the posterior decision boundary itself evolves ($P_{i,t}(Y|X) \neq P_{i,t+\Delta t}(Y|X)$). 

The confluence of Spatial Concept Skew and Temporal Concept Shift creates a highly complex problem space. To systematically categorize and evaluate existing solutions addressing these intersecting heterogeneities, this survey proposes a multi-dimensional taxonomy encompassing four key dimensions: Data Characteristics, Learning Architectures, Adaptation Strategies, and Application Domains.

\#\#\# 1. Data Characteristics and Drift Detection
Detecting and characterizing General Concept Drift is a prerequisite for any adaptive FL system. Historically, centralized drift detection methods such as the Drift Detection Method (DDM) and Early Drift Detection Method (EDDM) [CITE:101, CITE:102] relied on continuous, unhindered access to raw data streams to monitor error rate fluctuations and statistical deviations. 

However, transitioning these centralized paradigms to FL introduces severe architectural and regulatory constraints. Traditional drift detection assumes a global, unified view of the data stream. In contrast, FL environments are bound by strict privacy mandates that explicitly prohibit raw data sharing. Furthermore, communication bottlenecks and intermittent client availability restrict the high-frequency, centralized monitoring required by traditional DDM and EDDM. Consequently, a direct application of centralized drift detection to FL is fundamentally unfeasible.

Addressing these constraints, the literature has pivoted toward decentralized, privacy-preserving detection methodologies [CITE:328, CITE:330]. Autoencoder architectures are repeatedly utilized for robust representation learning and anomaly detection under drifting conditions, as they can evaluate reconstruction errors locally without exposing raw data [CITE:195, CITE:349, CITE:607]. For instance, spatio-temporal memory-enhanced autoencoders have been successfully deployed for cloud-edge collaborative fault diagnosis \cite{Jiang2024}. Furthermore, advanced frameworks have integrated dual drift detectors using metrics like Rényi divergence for abrupt shifts and Wasserstein distance for gradual trends. By utilizing these lightweight statistical measures, systems can trigger selective client retraining without violating FL privacy or communication constraints \cite{Moayyed2025}.

\#\#\# 2. Learning Architectures
Traditional static, round-based FL is increasingly recognized as insufficient for continuous streaming data [CITE:195, CITE:325, CITE:330]. A notable debate within architectural design concerns model aggregation and synchronization under drift. Early literature frequently posited that asynchronous updates were strictly necessary to handle non-stationary streaming data, arguing that traditional synchronous parameter averaging introduces significant synchronization bottlenecks and straggler effects. Consequently, asynchronous architectures, such as the Asynchronous Online Federated Learning (ASO-Fed) framework \cite{Chen2020}, have been frequently advocated to handle real-time non-stationary environments effectively.

However, this absolute claim has been nuanced by recent studies. Synchronous FL frameworks equipped with robust straggler mitigation and proximal regularization (e.g., variants of FedProx or FedNova) can also maintain stability in drifting environments, provided that the drift velocity does not exceed the communication and synchronization frequency [CITE:888, CITE:889]. Thus, the choice between synchronous and asynchronous aggregation depends heavily on the ratio of drift rate to network latency, rather than a strict mandate for asynchronous designs. Additionally, Clustered Federated Learning (CFL) has emerged as a dominant architectural solution for intersecting heterogeneities, grouping participants with similar data distributions or drift trajectories to train specialized models [CITE:216, CITE:227, CITE:454], while personalized FL frameworks utilize fine-grained layer aggregation to maintain robust performance across heterogeneous clients [CITE:137, CITE:219].

\#\#\# 3. Adaptation Strategies
Once drift is detected within a chosen architecture, systems must employ robust adaptation strategies. Applied frameworks have increasingly integrated continuous learning paradigms \cite{Miao2025} and online sequential extreme learning machines (OS-ELM) \cite{Ito2021}. OS-ELM is particularly advantageous in this context because it provides lightweight, sequential updating capabilities—requiring only single-pass learning without iterative backpropagation—making it ideal for rapid adaptation to evolving data distributions.

Interestingly, there is a divergence in knowledge-sharing mechanisms during adaptation. While most frameworks share model parameters or gradients [CITE:349, CITE:592], some approaches explicitly avoid parameter sharing to mitigate label heterogeneity. Instead, they share alternative representations such as "diagnostic knowledge"—comprising model predictions, soft labels, or feature prototypes. This strategy reduces communication overhead and preserves privacy while effectively tracking drift across disparate local distributions \cite{Jiang2024}. Furthermore, advanced adaptation strategies have begun employing aggregated Shapley values to dynamically evaluate and weight client contributions. This ensures that clients experiencing severe Temporal Concept Shift do not corrupt the global model during aggregation, isolating obsolete concepts while integrating relevant updates [CITE:412, CITE:415].

\#\#\# 4. Application Domains
A prominent pattern in the literature is the heavy focus on edge computing and Internet of Things (IoT) environments [CITE:328, CITE:349, CITE:592, CITE:607]. In these application domains, devices are highly resource-constrained (in terms of compute, memory, and battery) yet face continuous, high-velocity data streams. The constraints of these domains dictate the necessity of the lightweight detection mechanisms and communication-efficient adaptation strategies discussed in the previous dimensions, highlighting the deeply interconnected nature of this taxonomy.

\#\#\# Critical Analysis and Open Challenges
Despite these advancements, a critical analysis of the literature reveals a significant technical disconnect between theoretical models and practical implementations. Theoretical frameworks, particularly those grounded in statistical mechanics \cite{Straat2018}, model continuous drift as a strict Markovian process. These models mathematically demonstrate that when the drift velocity exceeds a critical threshold relative to the learning rate, the global model inevitably falls into persistent, sub-optimal plateau states—termed "retarded learning phases"—due to the delayed propagation of new concepts across the decentralized network. 

Conversely, applied continuous learning frameworks frequently report successful mitigation of catastrophic forgetting and continuous performance improvements. This theory-practice gap exists because empirical systems employ practical workarounds that fundamentally alter the theoretical assumptions of pure online learning. Mechanisms such as global replay buffers populated with synthetic data \cite{Miao2025} or sequential memory updates \cite{Ito2021} artificially inject historical or pseudo-stationary data into the current training process. By doing so, they violate the strict Markovian assumption of the theoretical models, artificially reducing the temporal variance of the data stream. However, these empirical workarounds introduce new vulnerabilities; synthetic data generation heavily depends on the generative capabilities of the global model and is highly susceptible to mode collapse, bias amplification, and quality degradation over time. Ultimately, statistical learning theory dictates that the temporal evolution of the joint data distribution $P_{i,t}(X,Y)$ must be explicitly bounded to guarantee convergence. Ignoring this variance in theoretical bounds while artificially suppressing it in empirical setups leaves a profound gap in the literature, requiring new theoretical frameworks that accurately model pseudo-stationary interventions.

Finally, several critical research gaps remain. Theoretical analyses currently rely on asymptotic limits or strict statistical assumptions (e.g., infinite network width or perfect Gaussian distributions) that rarely hold true in complex, real-world federated data streams. Furthermore, multi-model and clustering approaches often rely on unrealistic assumptions, such as having oracle knowledge of the exact number of underlying concepts or drift states. Most importantly, there is a critical absence of standardized, large-scale streaming datasets specifically designed to benchmark temporal and distributed General Concept Drift in FL. Current evaluations rely heavily on fragmented, proprietary, or entirely synthetic datasets, severely limiting the reproducibility and comparative analysis of emerging algorithms.

\begin{figure}[!t]
\centering
\includegraphics[width=\columnwidth]{figures/top_cited}
\caption{Top cited papers in the survey corpus.}
\label{fig:top_cited}
\end{figure}

\section{Critical Analysis}
\label{sec:critical_analysis}

Here is the revised "Critical Analysis" section, addressing all reviewer feedback while maintaining rigorous academic standards and preserving the requested citation format.

***

Despite the rapid proliferation of methodologies addressing concept drift in Federated Learning (FL), a critical examination of the literature reveals several fundamental tensions, unrealistic assumptions, and methodological gaps that hinder the transition of these frameworks into robust, real-world deployments. 

At the core of these challenges is the dual-heterogeneity problem, which requires precise mathematical delineation to prevent the terminology collisions commonly observed in the literature. In a standard FL environment, spatial heterogeneity (non-IID data) dictates that the joint data distributions across clients differ: $P_i(X, Y) \neq P_j(X, Y)$ for clients $i \neq j$. This spatial divergence manifests as covariate skew $P_i(X) \neq P_j(X)$, label skew $P_i(Y) \neq P_j(Y)$, or *concept skew* $P_i(Y|X) \neq P_j(Y|X)$. Crucially, this spatial skew must be explicitly disambiguated from temporal non-stationarity. When temporal evolution is introduced, the distributions change over time $t$, resulting in general *concept drift*, strictly defined as a change in the joint distribution: $P_{i,t}(X, Y) \neq P_{i,t+1}(X, Y)$. This temporal drift can be further decomposed into temporal covariate shift $P_{i,t}(X) \neq P_{i,t+1}(X)$ and temporal *concept shift* $P_{i,t}(Y|X) \neq P_{i,t+1}(Y|X)$. By harmonizing these definitions, it becomes evident that under dual-heterogeneity, the global learning objective is a dynamically moving target: $\min_w \sum_{i=1}^N p_i \mathbb{E}_{(x,y) \sim P_{i,t}} [\ell(w; x,y)]$.

A primary tension in the literature exists between theoretical models of these learning dynamics and the empirical successes reported in applied FL systems. Theoretical analyses, frequently grounded in online learning regret bounds or statistical mechanics, warn that learning in dual-heterogeneous environments often leads to persistent, sub-optimal plateau states \cite{Straat2018}. Furthermore, existing theoretical guarantees typically hinge on restrictive assumptions, such as asymptotic sample sizes, stationary mixing times, sub-Gaussian loss functions, or strict convexity. These constraints result in conservative bounds that systematically fail to capture the dynamics of finite-sample, non-convex practical scenarios where deep neural networks are deployed. Conversely, applied frameworks frequently report continuous performance improvements and the successful mitigation of catastrophic forgetting [CITE:195, 349, 592]. 

This theory-practice gap exists primarily because empirical workarounds artificially induce localized stationarity, effectively bypassing the strict non-stationary constraints of theoretical models. For example, global replay buffers utilizing synthetic data generation \cite{Miao2025} approximate a stationary historical distribution $P_{1:t}(X, Y)$ over a sliding window, transforming a non-stationary stream into a pseudo-stationary batch problem. Similarly, lightweight sequential learning models enable rapid adaptation by minimizing the temporal footprint of the training phase, avoiding the pitfalls of iterative backpropagation over drifting data \cite{Ito2021}. While these broader mechanisms demonstrate empirical utility, they introduce their own vulnerabilities. Generative replay heavily depends on the global model's capacity, making it susceptible to bias and long-term quality degradation, while rapid sequential learners often sacrifice deep feature extraction capabilities. A unified theory explaining how these finite-memory empirical workarounds formally bypass asymptotic theoretical limitations remains a critical open challenge.

Another critical debate pertains to model aggregation and synchronization paradigms under continuous drift. Traditional FL relies on synchronous parameter averaging, assuming clients can compute and communicate updates in uniform rounds. In edge computing and IoT environments experiencing continuous streaming data, synchronous updates face severe bottlenecks due to stragglers and real-time processing constraints [CITE:328, 592]. Consequently, early assertions in the literature posited that asynchronous online federated learning frameworks are strictly necessary for handling non-stationary streams effectively [CITE:196, 592]. However, this absolute stance has been increasingly challenged. Synchronous FL frameworks, when augmented with robust straggler mitigation and proximal regularization (e.g., variants of FedProx or FedNova), demonstrate substantial resilience and remain highly viable in drifting environments, provided the drift rate does not outpace the aggregation frequency [CITE:610, 611]. 

Furthermore, there is an ongoing debate regarding what information should be shared across the network. While the majority of frameworks share raw model parameters or gradients [CITE:349, 592], this approach is highly vulnerable to negative transfer when clients experience heterogeneous local drifts. Alternative approaches explicitly avoid parameter sharing, opting instead to share "diagnostic knowledge"—such as localized drift detection statistics, loss landscape geometries, or abstract class prototypes [CITE:503, 514, 607]—to mitigate label heterogeneity. This divergence highlights a critical future opportunity to explore alternative knowledge-sharing mechanisms that can handle temporal drift and non-IID data simultaneously without compromising model plasticity.

The reliance on unrealistic assumptions also extends to the algorithmic design of adaptation strategies, particularly in multi-model and clustered federated learning (CFL) approaches [CITE:140, 216]. These frameworks group participants with similar $P_{i,t}(X, Y)$ distributions to train specialized models. However, they frequently rely on oracle knowledge, such as knowing the exact number of underlying concepts beforehand, or assuming the availability of balanced classifiers across clients. In dynamic, real-world environments where $P_{i,t}(X, Y)$ evolves unpredictably, these assumptions rarely hold. 

Compounding these algorithmic limitations is a significant methodological disconnect in drift detection. In traditional machine learning, centralized drift detection algorithms, such as the Drift Detection Method (DDM) or Early Drift Detection Method (EDDM), are considered the gold standard [CITE:112, 115]. These methods inherently assume full access to a continuous, centralized stream of data and immediate availability of true labels to monitor error rate fluctuations. However, transitioning these centralized paradigms directly to a federated context introduces severe operational constraints. In FL, data is strictly decentralized, privacy restrictions prohibit the pooling of client-level error statistics, and edge devices rarely have immediate access to ground-truth labels for incoming streams. Because centralized supervised methods are rendered impractical, federated concept drift detection must rely on unsupervised mechanisms. Unfortunately, current unsupervised federated detection methods often fall back on strict statistical assumptions—such as variables adhering to a Gaussian distribution or requiring perfectly stationary reference windows—which are easily violated in complex, high-dimensional federated data streams [CITE:420, 512].

Finally, the most pressing systemic gap in this domain is the critical absence of standardized, widely adopted benchmarking suites. Current evaluations of temporal and distributed concept drift in FL rely heavily on highly fragmented, domain-specific, or entirely synthetic datasets [CITE:584, 585]. Researchers frequently simulate non-stationarity by artificially partitioning static datasets using Dirichlet distributions [CITE:424, 584], an approach that fails to accurately reflect the complex, spatio-temporal evolution of real-world streaming data. Without unified, large-scale streaming datasets specifically designed to benchmark federated concept drift, cross-study empirical comparisons remain highly challenging, obscuring the true efficacy and scalability of proposed detection and adaptation algorithms. 

Resolving these overlapping challenges—ranging from theoretical gaps and synchronization debates to benchmarking deficiencies—necessitates a structured understanding of the problem space. To bridge the gap between theoretical models and practical deployments, it is imperative to systematically map the landscape of current solutions. This critical need directly motivates the multi-dimensional taxonomy proposed in the subsequent sections of this survey. By categorizing the existing literature across four key dimensions—Learning Architectures, Adaptation Strategies, Application Domains, and Data Characteristics—the following sections provide a comprehensive framework to navigate these challenges and guide future research toward more robust, assumption-free FL systems.

\section{Future Directions and Open Problems}
\label{sec:future_directions}

While significant progress has been made in addressing non-stationary environments within federated learning (FL), several critical research gaps remain. Building upon the multi-dimensional taxonomy proposed in this survey—which categorizes FL systems across Data Characteristics, Adaptation Strategies, Learning Architectures, and Application Domains—we identify key future directions. A fundamental prerequisite for exploring these directions is the rigorous mathematical disambiguation of non-stationary phenomena, which has historically suffered from terminology collisions in the literature. Future research must explicitly and consistently delineate between general temporal *Concept Drift* (a shift in the joint data distribution over time, $P_t(X, Y) \neq P_{t+1}(X, Y)$), temporal *Concept Shift* (changes in the conditional distribution, $P_t(Y|X) \neq P_{t+1}(Y|X)$), spatial *Concept Skew* (client-level conditional shifts where $P_i(Y|X) \neq P_j(Y|X)$ for clients $i$ and $j$), and general spatial data heterogeneity ($P_i(X, Y) \neq P_j(X, Y)$). The intersection of these distinct spatiotemporal dynamics with distributed, resource-constrained edge architectures presents unique challenges that require targeted investigation.

**Data Characteristics: Standardized Benchmarks and Real-World Datasets**
A critical barrier to advancing FL under concept drift is the absence of unified, large-scale, real-world streaming datasets specifically designed to benchmark coupled temporal and spatial shifts. Currently, evaluations rely heavily on highly fragmented, domain-specific, or entirely synthetic datasets. This lack of standardized benchmarking suites makes cross-study empirical comparisons highly challenging. Future research must prioritize the development of comprehensive, open-source benchmarks that accurately reflect the complexities of real-world non-stationary environments. These benchmarks should formally quantify varying degrees of label delay and spatiotemporal distribution shifts (e.g., by measuring the divergence $D_{KL}(P_t \parallel P_{t+1})$). Most importantly, they must explicitly decouple temporal concept shift from spatial concept skew, providing a rigorous foundation for evaluating whether an algorithm is adapting to genuine temporal changes or merely overfitting to local spatial heterogeneity.

**Adaptation Strategies: Bridging Centralized Detection, Theory, and Continual Learning**
Historically, drift detection has relied heavily on centralized statistical process control algorithms, such as the Drift Detection Method (DDM) and Early Drift Detection Method (EDDM) [CITE:112, CITE:113]. These traditional mechanisms assume continuous, centralized access to a single stream of error rates to monitor statistical deviations. However, directly porting these centralized mechanisms to FL environments is fundamentally flawed. In FL, data is distributed, private, and subject to severe spatial heterogeneity. Transmitting local error streams to a central server violates privacy constraints and incurs prohibitive communication overhead. Furthermore, spatial concept skew can easily trigger false-positive drift alerts if a global model is evaluated on a highly skewed local dataset. A critical future direction involves redesigning these detection mechanisms to operate under strict FL constraints. Researchers must develop localized, skew-aware drift detection algorithms that can differentiate between temporal concept shift and spatial heterogeneity, minimizing server-client communication while preserving privacy.

Furthermore, there remains a significant disconnect between theoretical models of concept drift in federated settings and practical implementations. Theoretical frameworks frequently predict persistent sub-optimal plateau states or catastrophic forgetting, assuming unconstrained optimization over continuously shifting distributions or strict Markovian drift transitions \cite{Straat2018}. In contrast, applied frameworks frequently report continuous performance improvements by artificially stabilizing the local empirical risk landscape. For example, global replay buffers anchor the optimization process by reintroducing historical data \cite{Miao2025}, fundamentally breaking the Markovian assumption. Similarly, lightweight sequential learning algorithms provide exact sequential updating without iterative backpropagation, enabling rapid adaptation on edge devices that defies traditional gradient-descent bounds \cite{Ito2021}. Future theoretical work must reconcile this disconnect by developing finite-sample generalization bounds that explicitly incorporate these memory-augmented dynamics and explain empirical successes under realistic, relaxed statistical assumptions (e.g., non-Gaussian noise, bounded memory constraints, and non-Markovian state transitions). 

Finally, as continual learning techniques are increasingly integrated into FL adaptation strategies, current frameworks relying on synthetic data generation for global replay buffers \cite{Miao2025} remain highly vulnerable to bias. Because these methods depend on the generative capabilities of the global model, temporal concept shift can compound generative errors over time, leading to semantic drift in the replayed samples. Future directions should explore drift-resilient memory management and generative replay strategies that maintain high fidelity without compromising strict client privacy constraints.

**Learning Architectures: Hybrid Synchronization and Alternative Knowledge Sharing**
Traditional federated learning relies heavily on synchronous parameter averaging. While earlier literature often posited that asynchronous updates were strictly necessary for non-stationary streaming data, recent advancements demonstrate that synchronous FL frameworks equipped with robust straggler mitigation techniques—such as proximal regularization \cite{paper882} or normalized updates \cite{paper883}—remain highly viable in drifting environments. However, synchronous methods can still face significant communication bottlenecks under continuous, high-velocity streaming data. Consequently, asynchronous online federated learning frameworks \cite{Chen2020} offer a complementary alternative, allowing for real-time, localized adaptations without waiting for global consensus. Future architectures should explore hybrid or adaptive synchronization strategies that balance the theoretical stability and convergence guarantees of synchronous aggregation with the low-latency responsiveness of asynchronous updates. 

Additionally, there is a promising opportunity to explore alternative knowledge-sharing mechanisms beyond standard parameter or gradient aggregation to simultaneously handle temporal concept shift and spatial concept skew. Researchers should investigate the communication of "diagnostic knowledge"—low-dimensional representations of local data distributions or decision boundaries, such as class-wise logits, feature distribution statistics, or prototype vectors. Sharing these representations has been shown to effectively mitigate negative transfer in heterogeneous, drifting environments by decoupling local feature extraction from global concept alignment \cite{Jiang2024}. 

**Application Domains: Robust Multi-Model and Dynamic Clustering**
To combat localized drift in complex application domains (e.g., IoT sensor networks and healthcare monitoring), many recent frameworks employ multi-model or clustered FL approaches \cite{Liu2025}. However, these methods often rely on restrictive assumptions, such as having oracle knowledge of the exact number of underlying concepts or requiring the communication of balanced classifiers, which are rarely available in dynamic, real-world environments. Future research must develop fully unsupervised, dynamic clustering mechanisms that do not require prior knowledge of the concept space and can organically grow or merge clusters as new spatiotemporal patterns emerge. Moreover, multi-model specialization approaches frequently suffer from a "cold-start" performance drop; when a new concept emerges, newly instantiated models exhibit poor initial performance due to insufficient training on the novel distribution $P_{t+1}(X, Y)$. Addressing this cold-start problem through rapid adaptation techniques, such as meta-learning or zero-shot transfer \cite{paper412}, is a vital direction for ensuring the reliable deployment of FL systems in high-stakes, non-stationary application domains.

\section{Conclusion}
\label{sec:conclusion}

This survey has presented a comprehensive review of Federated Learning (FL) under concept drift, a critical intersection for deploying collaborative machine learning in dynamic, real-world environments. Existing literature has predominantly treated FL data heterogeneity (spatial non-stationarity) and concept drift (temporal non-stationarity) as isolated phenomena. In contrast, this paper contributes a unified framework that explicitly addresses their intersection, providing a novel synthesis of how spatial and temporal dynamics compound one another. To structure this emerging field, we detailed a multi-dimensional taxonomy focused on the algorithmic and architectural mechanisms specific to FL drift adaptation. Moving beyond generic categorizations, our taxonomy structures the literature by Drift Detection Locus (distinguishing between server-centric and client-centric monitoring), Adaptation Mechanisms (encompassing model replacement, federated ensembles, and dynamic weighting), and Drift Synchronization across distributed nodes. At its core, the challenge of this domain lies in managing the spatiotemporal evolution of data, requiring systems to simultaneously navigate spatial concept skew across clients and temporal concept shift over time without relying on the continuous, unhindered data access assumed by traditional machine learning.

Addressing these compounded spatiotemporal challenges is significantly more complex in FL than in centralized paradigms. In traditional settings, concept drift detection algorithms rely on the continuous monitoring of raw data streams or sequential error rates. However, transitioning these monitoring techniques to a federated environment introduces fundamental operational conflicts: strict privacy requirements prohibit the central server from inspecting raw client data, while partial client participation and communication bottlenecks disrupt the sequential tracking required for timely drift detection. To overcome these constraints, the literature has pivoted toward FL-native countermeasures. As highlighted in our review, emerging solutions include personalized drift adaptation that tailors models to local temporal shifts, asynchronous update protocols that prevent lagging clients from bottlenecking adaptation, and federated ensemble methods that maintain diverse hypotheses to rapidly react to recurring drifts [CITE:195, 216, 330].

Our critical analysis of the synthesized literature reveals that a dominant theme in operationalizing these solutions is the necessity of transitioning from static, round-based FL to continuous, online learning paradigms tailored for streaming data [CITE:195, 325, 330, 592]. Edge computing and Internet of Things (IoT) environments frequently serve as the primary application domains, demanding communication-efficient and computationally lightweight adaptation mechanisms [CITE:328, 349, 592, 607]. To address the dual challenge of spatial skew and temporal shift, researchers have widely adopted clustered FL \cite{Liu2025} and robust representation learning utilizing autoencoder architectures [CITE:349, 607]. Despite these advancements, a profound divergence remains between theoretical models and practical implementations. Theoretical frameworks, often grounded in statistical mechanics and relying on strict asymptotic assumptions, warn of persistent sub-optimal plateau states induced by concept drift \cite{Straat2018}. Conversely, applied frameworks frequently bypass these theoretical limits, reporting continuous performance improvements through pragmatic, finite-sample workarounds such as generative replay buffers and adaptive sequential learning models [CITE:195, 349, 592]. This disconnect highlights the urgent need for a unified theory that bridges empirical successes with rigorous generalization bounds accommodating practical, finite-sample scenarios.

Moving forward, several critical research directions must be prioritized to mature this field. The most pressing empirical gap is the absence of standardized, large-scale, real-world streaming datasets specifically designed to benchmark distributed temporal concept drift, as current evaluations rely heavily on fragmented or entirely synthetic datasets. Furthermore, there is a significant opportunity to explore alternative knowledge-sharing mechanisms. Instead of standard gradient or parameter aggregation, sharing higher-level "diagnostic knowledge" represents a promising theoretical frontier to mitigate negative transfer in heterogeneous, drifting environments [CITE:503, 607]. While extracting this knowledge remains a theoretical bottleneck in many conventional systems, recent pioneering works demonstrate that transmitting localized, interpretable feature attributions or mathematical representations of decision boundary shifts (e.g., orthogonal projections of local gradients) allows the server to identify exactly *which* features are drifting without compromising privacy [CITE:503, 607, 801]. Additionally, researchers must address the "cold-start" performance degradation inherent in multi-model specialization approaches \cite{Jothimurugesan2022}, and improve the robustness of generative global replay buffers against bias and quality degradation over extended temporal horizons \cite{Miao2025}.

Ultimately, realizing the full potential of Federated Learning in the wild requires moving beyond the assumption of static data distributions. By embracing continuous, drift-aware collaborative learning frameworks—whether through optimized synchronous architectures, asynchronous updates, or dynamic federated ensembles—the research community can build resilient distributed artificial intelligence systems capable of adapting to the ever-evolving nature of real-world data streams.

\bibliographystyle{IEEEtran}
\bibliography{references}

\end{document}